\documentclass[11pt]{exam}
\usepackage[utf8]{inputenc}		% Caracteres latinos
\usepackage[spanish]{babel}		% Idioma español
\usepackage{geometry}			% Organizar el documento
\usepackage{graphicx}			% Incluir gráficos
\usepackage{makecell}			% Para personalizar las celdas de una tabla
\usepackage[nohdr]{mathexam}	% Añadimos el paquete mathexam (sin header)
\usepackage{amsmath}
\usepackage{amsfonts}
\usepackage{amssymb}
\usepackage{mathtools}
%\usepackage[]{mathptmx}        % A free version o Times Roman with mathematical symbols
%\usepackage{pzc}               % fuente cursiva (conjuntos) Zapf Chancery
%\usepackage{showframe}
%\usepackage{lipsum}

% DOCUMENTACIÓN DE LA CLASE EXAM
% http://ftp.inf.utfsm.cl/pub/tex-archive/macros/latex/contrib/exam/examdoc.pdf
% DOCUMENTACIÓN DE LA CLASE MATHEXAM
% http://ctan.dcc.uchile.cl/macros/latex/contrib/mathexam/doc/mathexam.pdf

% Definimos la geometría de la primera página
\geometry{
	a4paper,                    % Tamaño del documento
	hmargin = {1.5cm, 0.5cm}, 	% Margen horizontal izquierdo, derecho
	vmargin = {0.5cm, 1cm},	    % Margen vertical superior, inferior
	headsep = 4mm,				% Separación entre el encabezado y el texto
	head = 2.5cm,				% Tamaño del encabezado
	% marginparsep = 5mm, 		% Seperación entre las notas y el texto
	% marginpar = 1.5cm,		% Tamaño de las notas
	includeall,                 % incluye el encabezado, footer y notas dentro del tamaño del documento
	nomarginpar,	            % Elimina las notas
	foot = 1cm,                 % Tamaño del footer
	twoside,                	% Habilita el modo de impresión a doble cara
}

\selectlanguage{spanish}        % Selecciona el idioma

\pagestyle{headandfoot}         % Nuestro examen tendrá encabezado y pié

% DEFINIMOS EL ENCABEZADO
\header{
\begin{tabular}{l c c c l}
            \makecell{\includegraphics[height=2.5cm]{logo.png}} &
            \makecell{\textbf{IPEA 215} \\Raúl Scalabrini Ortiz} &
            \makecell{Examen} &
            \makecell{Curso\\1er Año} &
            \makecell[l]{Apellido y Nombre:\enspace\makebox[2in]{\hrulefill}\\Fecha: \today}
        \end{tabular}}{}{}

% DEFINIMOS EL PIE
\rfoot{Página \thepage\ de 
umpages}

% DOCUMENTO
\begin{document}

\ExamInstrBox{\centering Por favor, muestre todo el desarrollo de su trabajo! No se dará crédito a las respuestas que no tengan trabajo de apoyo. Escriba las respuestas en los espacios provistos. Tiene 1 hora y 20 minutos para completar este examen.}

\pointpoints{punto}{puntos}
\pointformat{\bfseries\boldmath(\thepoints)}
\begin{questions}
    \question[1] ¿Qué gusto tiene la sal?
    
    \question[10] Contestar las siguientes preguntas:
    \begin{parts}
        \part[5]    ¿Qué es una fracción equivalente?
        \part[5]    De un ejemplo de suma de fracciones.
    \end{parts}
    
    \question[30] Resolver los siguientes cálculos combinados:
    \begin{parts}
        \part[10]    $ -3 + \frac{1}{2} \cdot 5= $
        \part[10]    $ x^2-1=0 $
        \part[10]    $ a_n = \frac{1}{n} + 3 \cdot \frac{1}{n} $
    \end{parts}
\end{questions}

% Geometría para la otra carilla

ewgeometry{
	hmargin = {1.5cm, 0.5cm},
	vmargin = {0.5cm, 1cm},
	%nofoot,			% Elimina el pié
	nohead,			% Elimina el encabezado
	nomarginpar,	% Elimina las notas
	includeall,
}% \savegeometry{geometria_1}

\pagestyle{foot}    % El estilo de ésta página sólo constará de pié de página
\runningfooter{}{}{Página \thepage\ de 
umpages}


%\lipsum[1-5]

% \restoregeometry
% \loadgeometry{geometria_1}


\end{document}
